\section{Two-weight threshold slices: a deficit-paid covariance and one-heavy-input theorem}
All logarithms in the proof are natural. Let $v_i=4q_i(1-q_i)$,
$P_f(S)=\widehat f_\mu(S)^2$, $v_S=\prod_{i\in S}v_i$, and write
\[
 E(f)=\sum_{S\ne\varnothing}P_f(S)\ln(v_S/P_f(S)),\quad
 E_{\ge2}(f)=\sum_{|S|\ge2}P_f(S)\ln(v_S/P_f(S)),\quad
 \Delta(f)=I(f)-\operatorname{Var}(f).
\]
Take the product-biased symmetric-threshold baseline
\[
 E(g)\le c\Delta(g),\qquad
 c=C_T\ln2=460+80\ln2,
\]
from the 38-page paper, Theorem 6.6, for every symmetric Boolean
threshold $g$, all arities and biases. Constants have zero entropy and
deficit. No weighted FEI+ hypothesis is assumed.

\paragraph{A universal singleton observation.}
For any Boolean sign function, increasing or not, let
$\rho_i=\mathbb P(f(X^{i=1})\ne f(X^{i=0}))$ and
$a_i=\mathbb E[(f(X^{i=1})-f(X^{i=0}))/2]$. Then
$|a_i|\le\rho_i$, $I_i=v_i\rho_i$, and $P_f(\{i\})=v_i a_i^2$.
Consequently
\[
 P_f(\{i\})\ln\frac{v_i}{P_f(\{i\})}
 =2v_i a_i^2\ln(1/|a_i|)
 \le2v_i(|a_i|-a_i^2)\le2(I_i-P_f(\{i\})).
\]
The first inequality uses $-x\ln x\le1-x$ for $0\le x\le1$.
Since $\sum_i(I_i-P_f(\{i\}))\le2\Delta(f)$, singleton
corrected entropy is at most $4\Delta(f)$ in nats.

\paragraph{Disjoint pivotal boundaries.}
Let $f_0=2\mathbf1_{\{N\ge k_0\}}-1$ and
$f_1=2\mathbf1_{\{N\ge k_1\}}-1$ on the same independent biased
bits, $N=\sum_jX_j$, with distinct integer thresholds $k_1<k_0$.
Constants are allowed. Their discrete derivatives in coordinate $j$
are indicators of $N-X_j=k_0-1$ and $N-X_j=k_1-1$ and have disjoint
supports. The cross Dirichlet form is therefore exactly zero:
\begin{equation}\label{a041s2:cross}
 \sum_{S\ne\varnothing}|S|\widehat f_0(S)\widehat f_1(S)=0.
\end{equation}
Put $C=\operatorname{Cov}(f_0,f_1)$. Nestedness gives $C\ge0$;
if $\alpha_j=\mathbb P(f_j=1)$, then $C=4\alpha_0(1-\alpha_1)$.
Subtracting covariance from \eqref{a041s2:cross} and applying
Cauchy--Schwarz on the nonnegative weights $|S|-1$ gives
\begin{equation}\label{a041s2:covariance}
 C=-\sum_{|S|\ge2}(|S|-1)\widehat f_0(S)\widehat f_1(S)
 \le\sqrt{\Delta(f_0)\Delta(f_1)}.
\end{equation}
This is a deficit payment, not a physical conditional-TC payment.

\paragraph{An interval inherits a strong bound from its two tails.}
Let $d=\mathbf1_{\{k_1\le N<k_0\}}$ and $g=2d-1=f_1-f_0-1$.
Disjoint derivative boundaries imply $I(g)=I(f_0)+I(f_1)$.
Also $\operatorname{Var}(g)=V_0+V_1-2C$, so
\begin{equation}\label{a041s2:banddeficit}
 \Delta(g)=\Delta_0+\Delta_1+2C
 \le2(\Delta_0+\Delta_1).
\end{equation}
For $\Phi_s(x)=x\ln(s/x)$, the elementary coefficient inequality
\begin{equation}\label{a041s2:coherent}
 \Phi_s((a-b)^2)\le2\Phi_s(a^2)+2\Phi_s(b^2)+2a^2+2b^2
 \qquad(a^2,b^2\le s)
\end{equation}
is equation (62) of the 38-page paper. A direct proof: if
$2a^2+2b^2\le s/e$, use monotonicity of $\Phi_s$ up to $s/e$, then
its subadditivity, and $\Phi_s(2x)\le2\Phi_s(x)$; otherwise use
$\Phi_s(x)\le s/e<2a^2+2b^2$ and $\Phi_s(a^2),\Phi_s(b^2)\ge0$.
Apply this only on $|S|\ge2$. The overhead uses
$\sum_{|S|\ge2}P_{f_j}(S)\le\Delta_j$, not the full variance.
Thus
\[
 E_{\ge2}(g)\le(2c+2)(\Delta_0+\Delta_1),\qquad
 E(g)\le(2c+10)(\Delta_0+\Delta_1).
\]
The second inequality follows from the singleton bound and
\eqref{a041s2:banddeficit}. This also implies
$E(g)\le(2c+10)\Delta(g)$, but the stronger two-tail budget is the
one used below. Empty or constant intervals are interpreted directly.

\paragraph{The one-exceptional-weight theorem.}
Let $X\sim\operatorname{Ber}(p)$, independently of arbitrary biased
bits $Y_1,\ldots,Y_n$, and put
\[
 F=2\mathbf1_{\{wX+\sum_jY_j\ge t\}}-1,
 \qquad w>0,\ t\in\mathbb R.
\]
No restriction is imposed on $n,w,t,p$, or on the $Y$ biases.
Then
\begin{equation}\label{a041s2:oneheavy}
 E_{\ge2}(F)\le K\Delta(F),\quad K=4c+12+4/e,
 \qquad H^+_\mu(F)\le\frac{K+4}{\ln2}\Delta(F).
\end{equation}
In particular the constants are independent of the exceptional real
weight and all biases, and the class includes both recovered two-branch
weighted controls with $w=n$. The word ``exceptional'' is used because
the weight can also be below one.
The displayed corrected-entropy constant in bits is below $3000$:
$e>1+1+1/2+1/6=8/3$, and the four positive terms
$2\sum_{j=0}^3[(2j+1)3^{2j+1}]^{-1}$ in the series for $\ln2$
already exceed $3715/5360$. Hence
$320+(1856+4/e)/\ln2<320+(3715/2)/(3715/5360)=3000$.

Here is the full coherent calculation. The two $X$ restrictions are
$f_0$ and $f_1$ with effective integer thresholds $\lceil t\rceil$
and $\lceil t-w\rceil$. If they coincide, $X$ is irrelevant and the
baseline applies. Otherwise use the preceding two-tail calculation,
put $d=(f_1-f_0)/2$, $\rho=\mathbb E d$, and $v=4p(1-p)$.
For $B\subseteq[n]$ the exact coefficients are
\[
 \widehat F(B)=(1-p)\widehat f_0(B)+p\widehat f_1(B),\qquad
 \widehat F(\{X\}\cup B)=\sqrt v\,\widehat d(B).
\]
There is no incoherent replacement of the first expression. If
$D=(1-p)\Delta_0+p\Delta_1$ and $\delta=v\rho(1-\rho)$,
the law of total variance and physical influence averaging give
\begin{equation}\label{a041s2:ledger}
 \Delta(F)=D+\delta.
\end{equation}
For coefficients avoiding $X$ and having $|B|\ge2$, apply
\eqref{a041s2:coherent} with scaled amplitudes $(1-p)\widehat f_0(B)$
and $-p\widehat f_1(B)$. Since
\[
 \Phi_s(z^2a^2)=z^2\Phi_s(a^2)+2z^2a^2\ln(1/z),\quad
 z^2\ln(1/z)\le z/e,
\]
their total corrected entropy is at most
\[
 (2c+2+4/e)D.
\]
For coefficients containing $X$ and having nonempty $B$, use
$\widehat d(B)=\widehat g(B)/2$. Their total is exactly
\[
 \frac v4[E(g)+(\ln4)\operatorname{Var}(g)].
\]
The $E(g)$ term is at most $(2c+10)D$, since
$p(1-p)(\Delta_0+\Delta_1)\le D$. The variance term equals
$(\ln4)\delta$, because $\operatorname{Var}(g)=4\rho(1-\rho)$.
These two sectors are all of the high-degree spectrum of $F$, hence
\[
 E_{\ge2}(F)\le(4c+12+4/e)D+(\ln4)\delta
 \le K(D+\delta).
\]
Finally add the singleton bound $4\Delta(F)$ and divide by $\ln2$.
Constant restrictions cause no division by variance or deficit.
Deterministic inputs are frozen and removed, with the threshold adjusted;
zero weights are removed and negative weights are oriented.

\paragraph{An unbounded nonsymmetric prime family is covered.}
For every $n\ge3$, the gate
$F_n=2\mathbf1_{\{2X+Y_1+\cdots+Y_n\ge3\}}-1$ is essential and
prime under genuine disjoint one-bit substitution. Every light input
is essential on $X=1$ with all other light inputs zero; $X$ is essential
when exactly one light input is one. For any substituted block with
size at least two and a nonempty complement, a binary inner gate
would permit only two distinct cofactor rows on that complement.
For a block of $k\ge2$ light inputs with $l=n-k\ge1$ light inputs
outside it, block counts zero, one, and two give outside-score
thresholds three, two, and one. Scores one and two are both attained
outside (one light input or the exceptional input), so these rows
are distinct. If $l=0$, then $k=n\ge3$ and block counts zero, one,
and three give constant negative, the exceptional-input literal,
and constant positive rows. For a block consisting of $X$ and
$k\ge1$ light inputs, properness leaves $l\ge1$ light inputs outside.
Block scores zero, one, two, and three are attained. The resulting
thresholds three, two, one, and zero include at least the distinct
constant negative or upper tail, the threshold-one tail, and constant
positive rows. Thus every proper candidate block has at least three
rows. This proves primality by the necessary cofactor criterion,
without assuming that prime gates are a simpler universal class.
Equation \eqref{a041s2:oneheavy} consequently supplies a bias-uniform
constant on a nonsymmetric prime family of unbounded arity.

\paragraph{Exact interface for arbitrary numbers of exceptional inputs.}
For the full two-weight target
$F(A,B)=2\mathbf1_{\{w\sum_{i\in A}X_i+\sum_{j\in B}Y_j\ge t\}}-1$,
let $c_s(A)$ be the indicator of all heavy-count values producing the
same effective symmetric threshold $f_s(B)$; merge identical slices.
The cells form a partition, $\pi_s=\mathbb E c_s$, and
$F=\sum_s c_s f_s$ is an exact physical decomposition.
Distinct slices have disjoint pivotal boundaries on $B$, so the matrix
\[
 G_{ss}=\Delta(f_s),\qquad G_{su}=-\operatorname{Cov}(f_s,f_u)quad(s\ne u)
\]
is the Gram matrix of the vectors
$(\sqrt{|T|-1}\widehat f_s(T))_{|T|\ge2}$ and is positive semidefinite.
In particular
\begin{equation}\label{a041s2:matrix}
 2\sum_{s<u}\pi_s\pi_u\operatorname{Cov}(f_s,f_u)
 \le\sum_s\pi_s^2\Delta(f_s).
\end{equation}
For $m(B)=\mathbb E_A F=\sum_s\pi_sf_s(B)$, disjoint boundaries also give
$I_B(m)=\sum_s\pi_s^2I(f_s)$. Therefore the exact spectral budget is
\begin{equation}\label{a041s2:mixed}
 \sum_{S_A\ne\varnothing,S_B\ne\varnothing}|S_B|P_F(S_A,S_B)
 =\sum_s\pi_s(1-\pi_s)I(f_s)\le\Delta(F).
\end{equation}
Additionally
\[
 \Delta(F)=\sum_s\pi_s\Delta(f_s)+\delta_A,\quad
 \delta_A=I_A(F)-\operatorname{Var}_A(\mathbb E_BF)\ge0.
\]
The nonnegativity follows from the spectral expression, or Poincare
for the real-valued function $\mathbb E_BF$.
These are overlapping descriptions of one budget, not additional
reserves to add together: $\sum_s\pi_s\Delta(f_s)$ equals
$\sum_{S_B\ne\varnothing}(|S_B|-1)P_F(S_A,S_B)$, while
$\delta_A$ assigns $|S_A|P_F$ to each mixed set and
$(|S_A|-1)P_F$ to each pure-heavy set. Their sum assigns exactly
$(|S_A|+|S_B|-1)P_F$ once to every nonempty set.

\paragraph{Exact first unsupported line for the full target.}
The full coefficients are
$\widehat F(S_A,S_B)=\sum_s\widehat c_s(S_A)\widehat f_s(S_B)$,
a coherent sum. The matrix budget pays squared-energy covariance, not
logarithmic description of its complete Fourier set. It remains
unproved that the entropy of all mixed coefficients, together with the
pure-heavy higher-degree sector, is paid by the displayed nonreused
deficit ledger with one absolute constant when $|A|$ is arbitrary.
Replacing the coherent sum by a mixture, or summing a two-slice
inequality once for each heavy input, does not supply that result.
The one-exceptional-input theorem is unconditional relative to the
stated symmetric baseline; the all-two-weight target is still open.
